\documentclass[12pt,a4paper]{article}
\usepackage[utf8]{inputenc}
\usepackage{graphicx}
\usepackage[none]{hyphenat}
\usepackage{hyperref}
\usepackage{xcolor}

\definecolor{darkgreen}{RGB}{0, 150, 0}

\begin{document}
	\begin{titlepage}
		
		\begin{center}
			\includegraphics[width=0.5\textwidth]{QUT.jpg}\\
			[0.03\textheight]  
			\Large\textbf{Bachelor of IT (Computer Science)}\\
			\Large\textbf{Assignment 1}\\
			\large\textbf{CAB301 - Algorithms and Complexity}\\
			[0.02\textheight]
			\large\textsl{Dane Madsen}\\
			\large\textsl{n10983864@qut.edu.au}
		\end{center}
		
	\end{titlepage}
	\tableofcontents
	\newpage
	
	\section{IsValidId Method}
		\subsection{Algorithm Design}
			This method checks whether a provided job ID is valid. It achieves this 
			by checking that the provided ID is greater than the minimum valid ID (1) and 
			less than the maximum valid ID (999). If the ID is meets these criteria, the 
			method returns true indicating the ID is valid, otherwise it returns false 
			indicating the ID is invalid.\\

			\textbf{ALGORITHM} \textit{IsValidId(v)}\\
			\null\hspace{1cm}// Given an integer (\textit{v})\\
			\null\hspace{1cm}// Returns True if \textit{v} is a valid job ID\\
			\null\hspace{1cm}// Otherwise returns False\\
			\null\hspace{1cm}\textbf{if} \textit{$v \geq 1$} \textbf{and} \textit{$v \leq 999$}\\
			\null\hspace{2cm}\textbf{return} \textit{True}\\
			\null\hspace{1cm}\textbf{else}\\
			\null\hspace{2cm}\textbf{return} \textit{False}\\

		\subsection{Testing}
			This method uses unit testing to test that the method is correctly validating the 
			job ID. Using a for loop, the unit test first tests every valid ID (1-999) to ensure that 
			the method returns true. It then tests two invalid IDs (0 and 1000) to ensure that the method 
			returns false.\\

			\textbf{Test Result:} \textbf{\textcolor{darkgreen}{PASS}}\\

	\newpage

	\section{IsValidExecutionTime Method}
		\subsection{Algorithm Design}
			This method simply checks whether a provided job execution time is valid. It achieves 
			this by simply checking whether the execution time is greater than 0. If the execution 
			time is greater than 0, the method returns true indicating the execution time is valid,
			otherwise it returns false indicating the execution time is invalid.\\

			\textbf{ALGORITHM} \textit{IsValidExecutionTime(v)}\\
			\null\hspace{1cm}// Given an integer (\textit{v})\\
			\null\hspace{1cm}// Returns True if \textit{v} is a valid job execution time\\
			\null\hspace{1cm}// Otherwise returns False\\
			\null\hspace{1cm}\textbf{if} \textit{$v > 0$}\\
			\null\hspace{2cm}\textbf{return} \textit{True}\\
			\null\hspace{1cm}\textbf{else}\\
			\null\hspace{2cm}\textbf{return} \textit{False}

		\subsection{Testing}
			This method uses unit testing to test that the method is correctly validating the 
			job execution time. Using a for loop, the unit test first tests 100 valid execution times 
			(1-100) to ensure that the method returns true. It then tests one invalid execution time 
			(0) to ensure that the method returns false.\\

			\textbf{Test Result:} \textbf{\textcolor{darkgreen}{PASS}}\\

	\newpage

	\section{IsValidPriority Method}
		\subsection{Algorithm Design}
			This method checks whether a provided job priority is valid. It achieves this by 
			checking that the provided priority is greater than or equal to the minimum valid priority (1) and 
			less than or equal to the maximum valid priority (9). If the priority is meets these criteria, the 
			method returns true indicating the priority is valid, otherwise it returns false 
			indicating the priority is invalid.\\

			\textbf{ALGORITHM} \textit{IsValidPriority(v)}\\
			\null\hspace{1cm}// Given an integer (\textit{v})\\
			\null\hspace{1cm}// Returns True if \textit{v} is a valid job priority\\
			\null\hspace{1cm}// Otherwise returns False\\
			\null\hspace{1cm}\textbf{if} \textit{$v \geq 1$} \textbf{and} \textit{$v \leq 9$}\\
			\null\hspace{2cm}\textbf{return} \textit{True}\\
			\null\hspace{1cm}\textbf{else}\\
			\null\hspace{2cm}\textbf{return} \textit{False}\\
			
		\subsection{Testing}
			This method uses unit testing to test that the method is correctly validating the 
			job priority. Using a for loop, the unit test first tests every valid priority (1-9) to ensure that 
			the method returns true. It then tests two invalid priorities (0 and 10) to ensure that the method 
			returns false.\\

			\textbf{Test Result:} \textbf{\textcolor{darkgreen}{PASS}}\\

	\newpage

	\section{IsTimeReceived Method}
		\subsection{Algorithm Design}
			This method checks whether a provided job time received is valid. It achieves this by
			checking that the provided time received is greater than zero. If the time received is
			greater than zero, the method returns true indicating the time received is valid, otherwise
			it returns false indicating the time received is invalid.\\

			\textbf{ALGORITHM} \textit{IsTimeReceived(v)}\\
			\null\hspace{1cm}// Given a job time received (\textit{v})\\
			\null\hspace{1cm}// Returns True if \textit{v} is a valid time received\\
			\null\hspace{1cm}// Otherwise returns False\\
			\null\hspace{1cm}\textbf{if} \textit{$v > 0$}\\
			\null\hspace{2cm}\textbf{return} \textit{True}\\
			\null\hspace{1cm}\textbf{else}\\
			\null\hspace{2cm}\textbf{return} \textit{False}
		
		\subsection{Testing}
			This method uses unit testing to test that the method is correctly validating the 
			job time received. Using a for loop, the unit test first tests 100 valid time received 
			(1-100) to ensure that the method returns true. It then tests one invalid time received 
			(0) to ensure that the method returns false.\\

			\textbf{Test Result:} \textbf{\textcolor{darkgreen}{PASS}}\\

	\newpage
			
	\section{Add Method}
		\subsection{Algorithm Design}
			This method adds a job to the job collection. It achieves this by first checking that the job 
			doesn't already exist in the collection. If the job does already exist in the collection, the 
			method returns false indicating the job was not added to the collection. If the job does not 
			already exist in the collection, the method adds the job to the collection, increments the count 
			variable and returns true.\\

			\textbf{ALGORITHM} \textit{Add(v)}\\
			\null\hspace{1cm}// Let (\textit{n}) be count\\
			\null\hspace{1cm}// Given a job (\textit{v})\\
			\null\hspace{1cm}// Returns True if \textit{v} was added to the jobs array (\textit{J})\\
			\null\hspace{1cm}// Otherwise returns False\\
			\null\hspace{1cm}\textbf{for} \textit{$i \gets 0$} \textbf{in} \textit{n - 1} \textbf{do}\\
			\null\hspace{2cm}\textbf{if} \textit{$v.id = J[i].id$}\\
			\null\hspace{3cm}\textbf{return} \textit{False}\\
			\null\hspace{1cm}\textbf{else}\\
			\null\hspace{2cm}\textit{$J[n] \gets v$}\\
			\null\hspace{2cm}\textit{$n \gets n + 1$}\\
			\null\hspace{2cm}\textbf{return} \textit{True}\\

		\subsection{Testing}
			This method uses unit testing to test that the method is correctly adding jobs to the 
			collection. First, the unit test creates a dummy job with valid properties, and then creates 
			a dummy job collection. The unit test then attempts to add the dummy job to the dummy job 
			collection. If this first attempt is successful, the unit test then attempts to add the dummy 
			job to the dummy job collection again. If the second attempt is unsuccessful, the unit test 
			passes. This ensures that the method can add a new job, but will not overwrite a job if a job in 
			the job collection already has the same ID as the job being added.\\

			\textbf{Test Result:} \textbf{\textcolor{darkgreen}{PASS}}\\

	\newpage

	\section{Contains Method}
		\subsection{Algorithm Design}
			This method is used to check if a job exists in the job collection. It achieves this by 
			checking if there is a job in the collection with the same ID as the provided job ID. If there is 
			a job with the same ID the method returns true, otherwise it returns false.\\

			\textbf{ALGORITHM} \textit{Contains(v)}\\
			\null\hspace{1cm}// Let (\textit{n}) be count\\
			\null\hspace{1cm}// Given an integer (\textit{v})\\
			\null\hspace{1cm}// Returns True if a job with the ID \textit{v} exists in the jobs array (\textit{J})\\
			\null\hspace{1cm}// Otherwise returns False\\
			\null\hspace{1cm}\textbf{for} \textit{$i \gets 0$} \textbf{in} \textit{n - 1} \textbf{do}\\
			\null\hspace{2cm}\textbf{if} \textit{$v = J[i].id$}\\
			\null\hspace{3cm}\textbf{return} \textit{True}\\
			\null\hspace{1cm}\textbf{else}\\
			\null\hspace{2cm}\textbf{return} \textit{False}
			
		\subsection{Testing}


	\section{Find Method}
		\subsection{Algorithm Design}
			This method is used to find a job in the job collection. It achieves this by using the provided 
			job ID to find the job in the collection. If the job is found, the method returns the job, 
			otherwise it returns null.\\

			\textbf{ALGORITHM} \textit{Find(v)}\\
			\null\hspace{1cm}// Let (\textit{n}) be count\\
			\null\hspace{1cm}// Given an integer (\textit{v})\\
			\null\hspace{1cm}// Returns the job with the ID \textit{v} if it exists in the jobs array (\textit{J})\\
			\null\hspace{1cm}// Otherwise returns null\\
			\null\hspace{1cm}\textbf{for} \textit{$i \gets 0$} \textbf{in} \textit{n - 1} \textbf{do}\\
			\null\hspace{2cm}\textbf{if} \textit{$v = J[i].id$}\\
			\null\hspace{3cm}\textbf{return} \textit{J[i]}\\
			\null\hspace{1cm}\textbf{return} \textit{null}\\

		\subsection{Testing}

	\section{Remove Method}
		\subsection{Algorithm Design}
			This method is used to remove a job from the job collection. It achieves this by using the 
			provided job ID to find the job in the collection. If the job is found, the method removes the 
			job from the collection, decrements the count variable and returns true. If the job is not found, 
			the method returns false.\\

			\textbf{ALGORITHM} \textit{Remove(v)}\\
			\null\hspace{1cm}// Let (\textit{n}) be count\\
			\null\hspace{1cm}// Given an integer (\textit{v})\\
			\null\hspace{1cm}// Returns True if a job with the ID \textit{v} was removed\\
			\null\hspace{1cm}// from the jobs array (\textit{J})\\
			\null\hspace{1cm}// Otherwise returns False\\
			\null\hspace{1cm}\textbf{for} \textit{$i \gets 0$} \textbf{in} \textit{n - 1} \textbf{do}\\
			\null\hspace{2cm}\textbf{if} \textit{$v = J[i].id$}\\
			\null\hspace{4cm}\textbf{for} \textit{$j \gets 0$} \textbf{in} \textit{n - 2} \textbf{do}\\
			\null\hspace{5cm}\textit{$J[j] \gets J[j + 1]$}\\
			\null\hspace{4cm}\textit{$n \gets n - 1$}\\
			\null\hspace{4cm}\textbf{return} \textit{True}\\
			\null\hspace{1cm}\textbf{return} \textit{False}\\
			
		\subsection{Testing}

			
	\section{ToArray Method}
		\subsection{Algorithm Design}
			This method is used to convert the job collection to an array. It achieves this by creating a 
			new array of the same size as the job collection and then copying the jobs from the job 
			collection to the new array. The method then returns the new array.\\

			\textbf{ALGORITHM} \textit{ToArray()}\\
			\null\hspace{1cm}// Let (\textit{n}) be count\\
			\null\hspace{1cm}// Returns a new array of copied from the jobs array (\textit{J})\\
			\null\hspace{1cm}\textit{$A \gets new Job[n]$}\\
			\null\hspace{1cm}\textbf{for} \textit{$i \gets 0$} \textbf{in} \textit{n - 1} \textbf{do}\\
			\null\hspace{2cm}\textit{$A[i] \gets J[i]$}\\
			\null\hspace{1cm}\textbf{return} \textit{A}\\
		
		\subsection{Testing}
	
	\section{FirstComeFirstServed Method}
		\subsection{Algorithm Design}
			This method is used to sort jobs for first come first served scheduling. It achieves 
			this by using the selection sort algorithm to sort the jobs by their arrival time.\\

			\textbf{ALGORITHM} \textit{FirstComeFirstServed()}\\
			\null\hspace{1cm}// Returns a new array of jobs sorted by their arrival time\\
			\null\hspace{1cm}\textit{$A \gets Jobs.ToArray()$}\\
			\null\hspace{1cm}\textbf{for} \textit{$i \gets 0$} \textbf{in} \textit{A.Length - 1} \textbf{do}\\
			\null\hspace{2cm}\textbf{for} \textit{$j \gets i + 1$} \textbf{in} \textit{A.Length} \textbf{do}\\
			\null\hspace{3cm}\textbf{if} \textit{$A[i].TimeRecieved > A[j].TimeRecieved$}\\
			\null\hspace{4cm}\textit{$temp \gets A[i]$}\\
			\null\hspace{4cm}\textit{$A[i] \gets A[j]$}\\
			\null\hspace{4cm}\textit{$A[j] \gets temp$}\\
			\null\hspace{1cm}\textbf{return} \textit{A}\\
		
		\subsection{Analysis}
			To implement selection sort, this method uses a nested for loop to iterate through the
			array of jobs. The outer loop performs n - 1 iterations, and the inner loop performs 
			1, 2, ..., n - 2, n - 1 iterations for every iteration of the outer loop, where n is 
			the length of the jobs array. Therefore, the total number of iterations can be expressed as 
			1 + 2 + ... + (n - 2) + (n - 1), and using the rule of sum of an arithmetic sequence 
			this equation can be further simplified to n(n - 1)/2. By dropping the constant factor and 
			the lower order terms, the total number of iterations can be expressed as O($n ^ 2$). Therefore, 
			the worst case time complexity of this method is O($n ^ 2$).\\
			
		\subsection{Testing}
	\newpage
	\section{Priority Method}
		\subsection{Algorithm Design}
			This method is used to sort jobs for priority scheduling. It achieves this by using the 
			insertion sort algorithm to sort the jobs by their priority.\\
			
			\textbf{ALGORITHM} \textit{Priority()}\\
			\null\hspace{1cm}// Returns a new array of jobs sorted by their priority\\
			\null\hspace{1cm}\textit{$A \gets Jobs.ToArray()$}\\
			\null\hspace{1cm}\textbf{for} \textit{$i \gets 1$} \textbf{in} \textit{A.Length} \textbf{do} \textit{$i++$}\\
			\null\hspace{2cm}\textbf{for} \textit{$j \gets i$} \textbf{in} \textit{0} \textbf{do} \textit{$j--$}\\
			\null\hspace{3cm}\textbf{if} \textit{$A[j].Priority < A[j - 1].Priority$}\\
			\null\hspace{4cm}\textit{$temp \gets A[j]$}\\
			\null\hspace{4cm}\textit{$A[j] \gets A[j - 1]$}\\
			\null\hspace{4cm}\textit{$A[j - 1] \gets temp$}\\
			\null\hspace{3cm}\textbf{else}\\
			\null\hspace{4cm}\textbf{break}\\
			\null\hspace{1cm}\textbf{return} \textit{A}\\
		
		\subsection{Analysis}
			To implement insertion sort, like the previous method this method uses a nested for loop to 
			iterate through the array of jobs. However, this time the inner loop iterates in reverse. 
			The outer loop performs n iterations, and the inner loop performs n - 1, n - 2, ..., 2, 1 
			iterations for every iteration of the outer loop, where n is the length of the jobs array. Therefore, 
			the total number of iterations can be expressed as (n - 1) + (n - 2) + ... + 2 + 1, and by again 
			using the rule of sum of an arithmetic sequence this equation can be further simplified to n(n - 1)/2. 
			This is the same as the previous method so therefore, the worst case time complexity of this method 
			is O($n ^ 2$).\\

		\subsection{Testing}
	\section{ShortestJobFirst Method}
		\subsection{Algorithm Design}
			This method is used to sort jobs for shortest job first scheduling. It achieves this by 
			using the bubble sort algorithm to sort the jobs by their execution time.\\
			
			\textbf{ALGORITHM} \textit{ShortestJobFirst()}\\
			\null\hspace{1cm}// Returns a new array of jobs sorted by their execution time\\
			\null\hspace{1cm}\textit{$A \gets Jobs.ToArray()$}\\
			\null\hspace{1cm}\textbf{for} \textit{$i \gets 0$} \textbf{in} \textit{$A.Length - 1$} \textbf{do}\\
			\null\hspace{2cm}\textbf{for} \textit{$j \gets 0$} \textbf{in} \textit{$A.Length - i - 1$} \textbf{do}\\
			\null\hspace{3cm}\textbf{if} \textit{$A[j].ExecutionTime > A[j + 1].ExecutionTime$}\\
			\null\hspace{4cm}\textit{$temp \gets A[j]$}\\
			\null\hspace{4cm}\textit{$A[j] \gets A[j + 1]$}\\
			\null\hspace{4cm}\textit{$A[j + 1] \gets temp$}\\
			\null\hspace{1cm}\textbf{return} \textit{A}\\

		\subsection{Analysis}
			To implement bubble sort, again this method uses a nested for loop to iterate through the array of 
			jobs. The outer loop performs n - 1 iterations, and the inner loop performs 1, 2, ..., n - 3, n - 2 
			iterations for every iteration of the outer loop, where n is the length of the jobs array. Therefore, 
			the total number of iterations can be expressed as 1 + 2 + ... + (n - 3) + (n - 2), and by again 
			using the rule of sum of an arithmetic sequence this equation can be further simplified to n(n - 1)/2. 
			This is the same as the previous two methods so therefore, the worst case time complexity of this 
			method is O($n ^ 2$).\\

		\subsection{Testing}

\end{document}